\scanpage{209}
sch de \uline{classes de chemins}, ils sont déterminés \inserted{(s'il y en a au moins un)} par les automorphismes des foncteurs, qui font les pts d'un groupe algébrique.
Mais en général il peut y avoir des pts à valeurs dans un même $k$ non isomorphes.
Le groupe de Galois opérant sur \deleted{les points} l'ens des pts à valeurs dans $k$.
Si on a un pt $\xi$ à val. dans $k$, tous les autres pts s'en déduisent par des fibrés principaux \unsure{hom.} \unsure{du|au} (pro)groupe algébrique $\underline{\operatorname{Aut}}(\xi)$.

\diagram{\[
\left\{\begin{array}{ll}
\text{Pts $\ell$-adiques}&\text{valeurs }\mathbf Z_\ell\ \unsure{\text{ou}|-}\ \mathbf Q_\ell\\[6pt]
\text{Pts transcendants}&\text{valeurs }\mathbf R\ \deleted{\text{ou }\mathbf C\ \inserted{\unsure{\text{alg.}|\text{sup.}}}}\ \text{voire }\mathbf Z\\
&\begin{gathered}
\text{[structure supplémentaire]}\\
\text{Je doute que ces pts correspondent}\\
\text{biunivoquement aux pts de}\\
X\text{ à valeurs dans }\mathbf C,\\
\text{mod conjugaison}
\end{gathered}
\end{array}\right.
\]
\[
\begin{array}{ll}
\text{Pts Hodge-de Rham}&\text{valeur }k\text{ quelconque}\\
\text{pts de Witt}&\text{valeurs }W_p(k)\quad\text{$k$ corps de car $p$}
\end{array}
\]}

\margin{208}
